% الجزء: نظرية البرهان
% الفصل: حذف القطع
% القسم: ce-topmost

\documentclass[../../../include/open-logic-section]{subfiles}

\begin{document}

\olfileid{pt}{cut}{top}

\olsection{إزالة القطوع العليا}

\begin{defn}
تأمل !!a{proof}~$\pi$ ينتهي بقاعدة \CutCS{}،
\[
\AxiomC{}
\RightLabel{$\pi_1$}
\Deduce$\Gamma \fCenter \Delta, !A$
\AxiomC{}
\RightLabel{$\pi_2$}
\Deduce$!A, \Gamma \fCenter \Delta$
\RightLabel{\CutCS}
\BinaryInf$\Gamma \fCenter \Delta$
\DisplayProof
\]
ولا يحتوي على استدلالات \CutCS{} أخرى. يكون \emph{ارتفاع القطع} في~$\pi$
هو $\cheight{\pi}=\pheight{\pi_1}+\pheight{\pi_2}$. وتكون
\emph{رتبة القطع} في~$\pi$ هي $\cutrank{\pi}=\depth{!A}$.
\end{defn}

\begin{lem}\ollabel{lem:cut-adm-G3c}
قاعدة \CutCS{} مقبولة في~\Log{G3c}. وبعبارة أخرى، إذا كان $\pi$ هو
!!a{proof} ينتهي بــ\CutCS{} واحدة، ولا يستعمل فيما عدا ذلك إلا قواعد
النسق~$\Log{G3c}$، فهناك !!a{proof}~$\pi'$ للمتتالية النهائية نفسها في~\Log{G3c}.
\end{lem}

\begin{proof}
البرهان باستقراء مزدوج على !!{height} القطع ورتبته. ينتهي
!!{proof}~$\pi$ بما يأتي:
\[
\AxiomC{}
\RightLabel{$\pi_1$}
\Deduce$\Gamma \fCenter \Delta, !A$
\AxiomC{}
\RightLabel{$\pi_2$}
\Deduce$!A, \Gamma \fCenter \Delta$
\RightLabel{\CutCS}
\BinaryInf$\Gamma \fCenter \Delta$
\DisplayProof
\]

أساس الاستقراء: $\cheight{\pi}=0$ و$\cutrank{\pi}=0$. لما كان
$\cheight{\pi}=\pheight{\pi_1}+\pheight{\pi_2}=0$، فإن كلًا من
$\Gamma \Sequent \Delta, !A$ و$!A, \Gamma \Sequent \Delta$ مسلّمة.
ولدينا حالتان: (أ)~إما أن $!A$ ليست رئيسة في إحداهما، أو (ب)~أن $!A$ رئيسة
فيهما كلتيهما. وسنبيّن أدناه في الحالتين أن $\Gamma \Sequent \Delta$ لا بد
أن تكون مسلّمة (ولا نحتاج إلى فرض الاستقراء لإثبات ذلك).

سنميّز الآن عددًا من الحالات الممكنة بحسب ما إذا كان $\pi_1$ و$\pi_2$
مسلّمتين أو ينتهيان باستدلال، وبحسب ما إذا كانت $!A$ رئيسة في ذلك الاستدلال
أم لا. تغطي الحالتان A وB أساس الاستقراء. وفي خطوة الاستقراء نفترض أن
$\cheight{\pi}>0$ أو $\cutrank{\pi}>0$ (أو كليهما). ويمكننا استعمال فرض
الاستقراء، أي إن \CutCS{} الأخيرة يمكن حذفها من أي !!{proof} ينتهي
بــ\CutCS{} واحدة، وتكون إما رتبة قطعه $=\cutrank{\pi}$ و!!{height} قطعه
$<\cheight{\pi}$، وإما رتبة قطعه $<\cutrank{\pi}$. والحالة التي يكون فيها
$\cheight{\pi}=0$ (أي المقدمتان مسلّمتان) تغطيها الحالتان A وB، والحالات
التي يكون فيها $\cheight{\pi}>0$ تغطيها الحالات C--D.

\paragraph{الحالة A:}
ترد $!B$ ما في كل من $\Gamma$ و~$\Delta$.
عندئذ تكون $\Gamma \Sequent \Delta$ مسلّمة (إذا كانت $!B$ ذرية)، أو يكون
لها !!a{proof}~$\pi'$ في \Log{G3c} من دون \CutCS{} بحسب
\olref[seq][ex]{prop:G3c-axioms}. وهذا يغطي الحالة~(أ) من أساس الاستقراء،
لأنه إذا لم تكن $!A$ رئيسة في $\Gamma \Sequent \Delta, !A$ أو لم تكن
رئيسة في $!A, \Gamma \Sequent \Delta$، وكانت هاتان مسلّمتين، فلا بد أن
ترد $!B$ ذرية أخرى في كل من $\Gamma$ و~$\Delta$. كما يغطي خطوة الاستقراء
إذا كانت إحدى المقدمتين مسلّمة لا تكون فيها $!A$ رئيسة (حتى إن كانت
المقدمة الأخرى نتيجة قاعدة).

\paragraph{الحالة B:} المقدمتان مسلّمتان، و$!A$ رئيسة، ومن ثم ذرية. تأمل
المقدمة اليسرى $\Gamma \Sequent \Delta, !A$. لما كانت $!A$ رئيسة، فلا بد
أن ترد في $\Gamma$ (وإلا لما كانت $\Gamma \Sequent \Delta, !A$ مسلّمة).
وبالمثل، لا بد أن ترد $!A$ في $\Delta$، وإلا لما كانت $!A, \Gamma
\Sequent \Delta$ مسلّمة. وهكذا تكون $!A$ ذرية وترد في كل من $\Gamma$
و$\Delta$، أي إن $\Gamma \Sequent \Delta$ مسلّمة. وهذا يغطي الحالة~(ب)
من أساس الاستقراء.

\begin{center}
عرض بديل للحالتين A وB
\end{center}

\paragraph{الحالة A:} إحدى مقدمتي القطع مسلّمة من الصورة $!C, \Gamma'
\Sequent \Delta', !C$، حيث $!C$ ذرية. إذا لم تكن $!C$ هي !!{formula}
القطع~$!A$، فلا بد أن ترد $!C$ في كل من $\Gamma$ و~$\Delta$. وعندئذ تكون
$\Gamma \Sequent \Delta$ مسلّمة، ويمكننا اتخاذها~$\pi'$. وإذا كانت
$!C\ident !A$ هي !!{formula} القطع، فإما أن يكون (إذا كانت المقدمة اليسرى
مسلّمة) $!A\in\Gamma$ و$\Gamma=\Gamma', !A$، أو يكون (إذا كانت المقدمة
اليمنى مسلّمة) $!A\in\Delta$ و$\Delta=\Delta', !A$. في الحالة الأولى،
يكون $\pi_2$ هو !!a{proof} لــ$!A, !A, \Gamma' \Sequent \Delta$، وبمقبولية
\LeftR{\Contraction} (\olref[seq][inv]{prop:G3c-cont-adm}) يوجد !!a{proof}
لــ$\Gamma \Sequent \Delta$. وفي الحالة الثانية، يكون $\pi_1$ هو
!!a{proof} لــ$\Gamma \Sequent \Delta', !A, !A$، وبمقبولية
\RightR{\Contraction} نحصل على !!a{proof} لــ$\Gamma \Sequent \Delta$.

\paragraph{الحالة B:} إحدى المقدمتين مسلّمة من الصورة~$\lfalse,
\Gamma' \Sequent \Delta'$. فإذا كانت المقدمة اليسرى مسلّمة من هذه الصورة،
كان $\lfalse\in\Gamma$، وكانت $\Gamma \Sequent \Delta$ مسلّمة أيضًا.
وإذا كانت المقدمة اليمنى مسلّمة من هذه الصورة، فإما أن يكون
$\lfalse\in\Gamma$ (فتكون $\Gamma \Sequent \Delta$ مسلّمة مرة أخرى)، أو
أن يكون $!A\ident\lfalse$. وفي الحالة الأخيرة، يكون $\pi_1$ هو !!a{proof}
لــ$\Gamma \Sequent \Delta, \lfalse$. ويمكننا تحويله إلى !!a{proof} لــ
$\Gamma \Sequent \Delta$ بحذف نسخة واحدة من~$\lfalse$ من تالي كل متتالية
في~$\pi_1$. (تمرين: تحقق من ذلك بالاستقراء على~$\pheight{\pi_1}$.)

افترض الآن أن الحالة~A ولا الحالة~B تنطبق. في الحالات الباقية يكون
$\cheight{\pi}>0$، أي إن مقدمة واحدة على الأقل هي نتيجة استدلال.

\paragraph{الحالتان C وD: تبديل ترتيب القطوع.}
تشترك الإمكانات المدروسة في الحالتين~C وD في نمط واحد. نبدأ بــ!!a{proof}
ينتهي بــ\CutCS{}، وفيه تُستنتج إحدى المقدمتين بقاعدة~$R$ لا تكون فيها
!!{formula} القطع~$!A$ رئيسة (بل تكون !!{formula} أخرى~$!D$ رئيسة). في
هذا !!{proof} تأتي \CutCS{} بعد القاعدة~$R$، ويكون شكله التخطيطي كما يأتي:
\[
\AxiomC{}
\RightLabel{$\pi_1$}
\Deduce$\Gamma \fCenter \Delta', !D, !A$
\AxiomC{}
\RightLabel{$\pi_3$}
\Deduce$!A, \Gamma, \Pi \fCenter \Delta', \Lambda$
\RightLabel{$R$}
\UnaryInf$!A, \Gamma \fCenter \Delta', !D$
\RightSubproofLabel{$\pi_2$}
\RightLabel{\CutCS}
\BinaryInf$\Gamma \fCenter \Delta', !D$
\DisplayProof
\]
لا يغطي هذا إلا الحالة التي تكون فيها $R$ قاعدة يمين ذات مقدمة واحدة، ولا
تكون $A$ هي !!{formula} القاعدة~$R$ الرئيسة، وترد !!{formula}~$!D$ التي
\emph{هي} رئيسة في تالي المقدمة اليمنى لــ\CutCS. أما إذا وردت $!D$ في
المقدم (أي كانت $R$ قاعدة يسار)، أو وردت في المقدمة اليسرى لــ\CutCS،
فسيختلف الشكل بالطبع. وتمثل !!{formula}s في $\Pi$ و~$\Lambda$
!!{formula}s الفعالة للقاعدة~$R$.

نعكس هذا الترتيب ونتأمل !!a{proof} تأتي فيه $R$ بعد~\CutCS{}:
\[
\AxiomC{}
\RightLabel{$\pi_1'$}
\Deduce$\Gamma, \Pi \fCenter \Delta', !D, \Lambda, !A$
\AxiomC{}
\RightLabel{$\pi_3'$}
\Deduce$!A, \Gamma, \Pi \fCenter \Delta', !D, \Lambda$
\RightLabel{\CutCS}
\BinaryInf$\Gamma, \Pi \fCenter \Delta', !D, \Lambda$
\RightLabel{$R$}
\UnaryInf$\Gamma \fCenter \Delta', !D, !D$
\DisplayProof
\]
نستعمل الإضعاف الحافظ للــ!!{height} لنحصل على $\pi_1'$ و$\pi_3'$ من
$\pi_1$ و$\pi_3$ على الترتيب.

ولأننا نبدّل ترتيب \CutCS{} و~$R$، يسمى ذلك «نقل القطع صعودًا عبر~$R$».
ويؤدي هذا إلى خفض ارتفاع !!{proof} المنتهي بالمقدمة اليمنى لــ\CutCS، ومن
ثم يخفض !!{height} القطع. أي يمكن افتراض أن !!{proof}s الفرعية المنتهية
بــ\CutCS{} الجديدة لا يزيد !!{height}ها على ارتفاع $\pi_1$ و$\pi_3$،
فتنخفض !!{height} القطع (لأننا حذفنا استدلالًا واحدًا من اليمين). وينطبق
الآن فرض الاستقراء، فيخبرنا بإمكان حذف استدلال \CutCS{}. وأخيرًا نحذف
النسخة الزائدة من~$!D$ باستعمال مقبولية~$\RightR{\Contraction}$.

\paragraph{الحالة C:}
ينتهي واحد على الأقل من $\pi_1$ و$\pi_2$ بقاعدة~$R$ ذات مقدمة واحدة لا
تكون فيها $!A$ رئيسة. وهناك 18 حالة في المجموع، بحسب ما إذا كانت $!A$ غير
رئيسة في المقدمة اليسرى أو اليمنى لــ\CutCS، وبحسب أي واحدة من القواعد
التسع ذات المقدمة الواحدة تكون~$R$ (\LeftR{\lnot}، \RightR{\lnot}،
\RightR{\lor}، \LeftR{\land}، \RightR{\lif}، وقواعد الكم الأربع). ونكتفي
بمثالين: أن تكون $!A$ غير رئيسة في الاستدلال الأخير لــ$\pi_2$، الذي
ينتهي إما بــ$\RightR{\lif}$ أو بــ\LeftR{\lexists}. أما الحالات الأخرى
فنظيرة لهما وتُترك تمارين.

افترض أن $\pi$ ينتهي بما يأتي:
\[
\AxiomC{}
\RightLabel{$\pi_1$}
\Deduce$\Gamma \fCenter \Delta', !B \lif !C, !A$
\AxiomC{}
\RightLabel{$\pi_3$}
\Deduce$!A, !B, \Gamma \fCenter \Delta', !C$
\RightLabel{$\RightR{\lif}$}
\UnaryInf$!A, \Gamma \fCenter \Delta', !B \lif !C$
\RightSubproofLabel{$\pi_2$}
\RightLabel{\CutCS}
\BinaryInf$\Gamma \fCenter \Delta', !B \lif !C$
\DisplayProof
\]
حيث $\Delta=\Delta', !B \lif !C$. نطبّق الإضعاف الحافظ للــ!!{height}
(\olref[seq][adm]{prop:weak-G3c-adm}) على $\pi_1$ لنحصل على !!a{proof}
$\pi_1'$ لــ$!B, \Gamma \fCenter \Delta', !B \lif !C, !C, !A$ (نضعف
بــ$!B$ على اليسار وبــ$!C$ على اليمين). وبالمثل نطبّق
\olref[seq][adm]{prop:weak-G3c-adm} على~$\pi_3$ لنحصل على
!!a{proof}~$\pi_3'$ لــ$!A, !B, \Gamma \fCenter \Delta', !B \lif !C,
!C$ (بالإضعاف بــ$!B \lif !C$ على اليمين). وينطبق فرض الاستقراء على
\[
\AxiomC{}
\RightLabel{$\pi_1'$}
\Deduce$!B, \Gamma \fCenter \Delta', !B \lif !C, !C, !A$
\AxiomC{}
\RightLabel{$\pi_3'$}
\Deduce$!A, !B, \Gamma \fCenter \Delta', !B \lif !C, !C$
\RightLabel{\CutCS}
\BinaryInf$!B, \Gamma \fCenter \Delta', !B \lif !C, !C$
\DisplayProof
\]
الذي !!{formula} قطعه~$!A$: فهذا !!{proof} له رتبة القطع نفسها، و!!{height}
قطع أقل من~$\pi$. وهذا يعطينا !!a{proof}~$\pi_4$ في \Log{G3c} من دون
\CutCS{} لــ$!B, \Gamma \fCenter \Delta', !B \lif !C, !C$. طبّق
$\RightR{\lif}$ لتحصل على !!a{proof} لــ$\Gamma \fCenter \Delta',
!B \lif !C, !B \lif !C$:
\[
\AxiomC{}
\RightLabel{$\pi_4$}
\Deduce$!B, \Gamma \fCenter \Delta', !B \lif !C, !C$
\RightLabel{$\RightR{\lif}$}
\UnaryInf$\Gamma \fCenter \Delta', !B \lif !C, !B \lif !C$
\DisplayProof
\]
وللحصول على !!a{proof}~$\pi'$ لــ$\Gamma \Sequent \Delta$، أي
$\Gamma \fCenter \Delta', !B \lif !C$، طبّق
\olref[seq][inv]{prop:G3c-cont-adm} (مقبولية الانكماش).

افترض الآن أن $\pi$ ينتهي بما يأتي:
\[
\AxiomC{}
\RightLabel{$\pi_1$}
\Deduce$\lexists[x][!B(x)], \Gamma' \fCenter \Delta, !A$
\AxiomC{}
\RightLabel{$\pi_3$}
\Deduce$!A, !B(c), \Gamma' \fCenter \Delta$
\RightLabel{$\LeftR{\lexists}$}
\UnaryInf$!A, \lexists[x][!B(x)], \Gamma' \fCenter \Delta$
\RightSubproofLabel{$\pi_2$}
\RightLabel{\CutCS}
\BinaryInf$\lexists[x][!B(x)], \Gamma' \fCenter \Delta$
\DisplayProof
\]
نطبّق فرض الاستقراء مرة أخرى على
\[
\AxiomC{}
\RightLabel{$\pi_1[!B(c) \Sequent {}]$}
\Deduce$\lexists[x][!B(x)], !B(c), \Gamma' \fCenter \Delta, !A$
\AxiomC{}
\LeftLabel{$\pi_3[\lexists[x][!B(x)] \Sequent {}]$}
\Deduce$!A, \lexists[x][!B(x)], !B(c), \Gamma' \fCenter \Delta$
\RightLabel{\CutCS}
\BinaryInf$\lexists[x][!B(x)], !B(c), \Gamma' \fCenter \Delta$
\RightLabel{\LeftR{\lexists}}
\UnaryInf$\lexists[x][!B(x)], \lexists[x][!B(x)], \Gamma' \fCenter \Delta$
\DisplayProof
\]
لاحظ أن شرط المتغير المميز متحقق، لأن شرط المتغير المميز على
\LeftR{\lexists} في~$\pi_2$ يضمن أن $c$ لا يرد في
$\lexists[x][!B(x)]$ أو $\Gamma'$ أو~$\Delta$. وأخيرًا، طبّق
\olref[seq][inv]{prop:G3c-cont-adm}.


\paragraph{الحالة D:}
ينتهي واحد على الأقل من $\pi_1$ و$\pi_2$ بقاعدة~$R$ ذات مقدمتين لا تكون
فيها $!A$ رئيسة. وهناك 6 حالات. وندرس الحالة التي لا تكون فيها $!A$ رئيسة
في الاستدلال الأخير لــ$\pi_2$، الذي ينتهي بــ$\RightR{\land}$؛ أما الحالات
الأخرى فنظيرة لها.

افترض أن $\pi$ ينتهي بما يأتي:
\[
\AxiomC{}
\RightLabel{$\pi_1$}
\Deduce$\Gamma \fCenter \Delta', !B \land !C, !A$
\AxiomC{}
\RightLabel{$\pi_3$}
\Deduce$!A, \Gamma \fCenter \Delta', !B$
\AxiomC{}
\RightLabel{$\pi_4$}
\Deduce$!A, \Gamma \fCenter \Delta', !C$
\RightLabel{$\RightR{\land}$}
\BinaryInf$!A, \Gamma \fCenter \Delta', !B \land !C$
\RightSubproofLabel{$\pi_2$}
\RightLabel{\CutCS}
\BinaryInf$\Gamma \fCenter \Delta$
\DisplayProof
\]
ينطبق فرض الاستقراء على !!{proof}s الآتية:
\[
\AxiomC{}
\RightLabel{$\pi_1'$}
\Deduce$\Gamma \fCenter \Delta', !B \land !C, !B, !A$
\AxiomC{}
\RightLabel{$\pi_3'$}
\Deduce$!A, \Gamma \fCenter \Delta', !B \land !C, !B$
\RightLabel{\CutCS}
\BinaryInf$\Gamma \fCenter \Delta', !B \land !C, !B$
\DisplayProof
\]
و
\[
\AxiomC{}
\RightLabel{$\pi_1''$}
\Deduce$\Gamma \fCenter \Delta', !B \land !C, !C, !A$
\AxiomC{}
\RightLabel{$\pi_4'$}
\Deduce$!A, \Gamma \fCenter \Delta', !B \land !C, !C$
\RightLabel{\CutCS}
\BinaryInf$\Gamma \fCenter \Delta', !B \land !C, !C$
\DisplayProof
\]
وهنا تُستخرج !!{proof}s~$\pi_1'$ و$\pi_1''$ في \Log{G3c} من~$\pi_1$
بالإضعاف الحافظ للــ!!{height} (\olref[seq][adm]{prop:weak-G3c-adm})
(مرة بالإضعاف بــ$!B$ على اليمين، ومرة بــ$!C$). وتُستخرج !!{proof}s
$\pi_3'$ و~$\pi_4'$ من $\pi_3$ و~$\pi_4$ أيضًا بالإضعاف الحافظ للــ
!!{height} من $\pi_3$ و~$\pi_4$ على الترتيب (بالإضعاف بــ$!B \land !C$
على اليمين).

يعطينا فرض الاستقراء !!{proof}s~$\pi_5$ و~$\pi_6$ في \Log{G3c} لكل من
$\Gamma \fCenter \Delta', !B \land !C, !B$ و$\Gamma \fCenter \Delta',
!B \land !C, !C$. ونجمعهما في !!a{proof} لــ$\Gamma \Sequent \Delta',
!B \land !C, !B \land !C$ باستعمال القاعدة~$\RightR{\land}$:
\[
\AxiomC{}
\RightLabel{$\pi_5$}
\Deduce$\Gamma \fCenter \Delta', !B \land !C, !B$
\AxiomC{}
\RightLabel{$\pi_6$}
\Deduce$\Gamma \fCenter \Delta', !B \land !C, !C$
\RightLabel{\RightR{\land}}
\BinaryInf$\Gamma \fCenter \Delta', !B \land !C, !B \land !C$
\DisplayProof
\]
وللحصول على !!a{proof} لــ$\Gamma \Sequent \Delta$، أي $\Gamma \fCenter
\Delta', !B \land !C$، طبّق \olref[seq][inv]{prop:G3c-cont-adm}
(مقبولية الانكماش).

\paragraph{الحالة E: اختزال القطوع.}
تتعلق الحالة الأخيرة بالوضع الذي ينتهي فيه كل من $\pi_1$ و$\pi_2$
باستدلال تكون فيه $!A$ رئيسة. وهناك ست حالات، واحدة لكل !!{main operator}
ممكن لــ$!A$.

في كل حالة، نحوّل \CutCS{} ذات !!{formula} القطع~$!A$ إلى قطع واحد أو أكثر
على !!{formula}s فرعية لــ$!A$، وهي تحديدًا !!{formula}s الفعالة في
الاستدلالات التي تكون $!A$ هي !!{formula}ها الرئيسة. ولذلك تسمى هذه
الحالات في الغالب «اختزال» استدلال القطع أو «تحويله».

\begin{enumerate}
\item إذا كان $!A \ident \lnot !B$، فإن $\pi$ ينتهي بــ
\[
\AxiomC{}
\RightLabel{$\pi_1'$}
\Deduce$!B, \Gamma \fCenter \Delta$
\RightLabel{\RightR{\lnot}}
\UnaryInf$\Gamma \fCenter \Delta, \lnot !B$
\LeftSubproofLabel{$\pi_1$}
\AxiomC{}
\RightLabel{$\pi_2'$}
\Deduce$\Gamma \fCenter \Delta, !B$
\RightLabel{\LeftR{\lnot}}
\UnaryInf$\lnot !B, \Gamma \fCenter \Delta$
\RightSubproofLabel{$\pi_2$}
\RightLabel{\CutCS}
\BinaryInf$\Gamma \fCenter \Delta$
\DisplayProof
\]
وبحسب فرض الاستقراء يمكن حذف \CutCS{} من
\[
\AxiomC{}
\RightLabel{$\pi_2'$}
\Deduce$\Gamma \fCenter \Delta, !B$
\AxiomC{}
\RightLabel{$\pi_1'$}
\Deduce$!B, \Gamma \fCenter \Delta$
\RightLabel{\CutCS}
\BinaryInf$\Gamma \fCenter \Delta$
\DisplayProof
\]

\item
إذا كان $!A \ident (!B \lor !C)$، فإن $\pi$ ينتهي بــ
\[
\AxiomC{}
\RightLabel{$\pi_1'$}
\Deduce$\Gamma \fCenter \Delta, !B, !C$
\RightLabel{\RightR{\lor}}
\UnaryInf$\Gamma \fCenter \Delta, !B \lor !C$
\LeftSubproofLabel{$\pi_1$}
\AxiomC{}
\RightLabel{$\pi_2'$}
\Deduce$!B, \Gamma \fCenter \Delta$
\AxiomC{}
\RightLabel{$\pi_2''$}
\Deduce$!C, \Gamma \fCenter \Delta$
\RightLabel{\LeftR{\lor}}
\BinaryInf$!B \lor !C, \Gamma \fCenter \Delta$
\RightSubproofLabel{$\pi_2$}
\RightLabel{\CutCS}
\BinaryInf$\Gamma \fCenter \Delta$
\DisplayProof
\]
تأمل !!{proof} المنتهي بــ\CutCS{}،
\[
\AxiomC{}
\RightLabel{$\pi_1'$}
\Deduce$\Gamma \fCenter \Delta, !B, !C$
\AxiomC{}
\RightLabel{$\pi_2'''$}
\Deduce$!C, \Gamma \fCenter \Delta, !B$
\RightLabel{\CutCS}
\BinaryInf$\Gamma \fCenter \Delta, !B$
\DisplayProof
\]
حيث يُستخرج $\pi_2'''$ من $\pi_2''$ بالإضعاف الحافظ للــ!!{height}.
رتبة قطعه $<\cutrank{\pi}$، لأن $\depth{!C}<\depth{!B \lor !C}$، ولذلك
يعطينا فرض الاستقراء !!a{proof} في \Log{G3c}، هو $\pi_1''$، لــ$\Gamma
\fCenter \Delta, !B$. أما !!{proof} المنتهي بــ\CutCS{}،
\[
\AxiomC{}
\RightLabel{$\pi_1''$}
\Deduce$\Gamma \fCenter \Delta, !B$
\AxiomC{}
\RightLabel{$\pi_2'$}
\Deduce$!B, \Gamma \fCenter \Delta$
\RightLabel{\CutCS}
\BinaryInf$\Gamma \fCenter \Delta$
\DisplayProof
\]
فرتبة قطعه $=\depth{!B}<\depth{!B \lor !C}$، ومن ثم ينطبق فرض الاستقراء
مرة أخرى، ويعطينا !!a{proof}~$\pi'$ لــ$\Gamma \fCenter \Delta$.

\item $!A \ident (!B \land !C)$. تمرين.

\item إذا كان $!A \ident (!B \lif !C)$، فإن $\pi$ ينتهي بــ
\[
\AxiomC{}
\RightLabel{$\pi_1'$}
\Deduce$!B, \Gamma \fCenter \Delta, !C$
\RightLabel{\RightR{\lif}}
\UnaryInf$\Gamma \fCenter \Delta, !B \lif !C$
\LeftSubproofLabel{$\pi_1$}
\AxiomC{}
\RightLabel{$\pi_2'$}
\Deduce$\Gamma \fCenter \Delta, !B$
\AxiomC{}
\RightLabel{$\pi_2''$}
\Deduce$!C, \Gamma \fCenter \Delta$
\RightLabel{\LeftR{\lif}}
\BinaryInf$!B \lif !C, \Gamma \fCenter \Delta$
\RightSubproofLabel{$\pi_2$}
\RightLabel{\CutCS}
\BinaryInf$\Gamma \fCenter \Delta$
\DisplayProof
\]
رتبة القطع في !!{proof} المنتهي بــ\CutCS{} الآتي أدنى من رتبتها في~$\pi$:
\[
\AxiomC{}
\RightLabel{$\pi_1'$}
\Deduce$!B, \Gamma \fCenter \Delta, !C$
\AxiomC{}
\RightLabel{$\pi_2''$}
\Deduce$!C, \Gamma \fCenter \Delta$
\RightLabel{\CutCS}
\BinaryInf$!B, \Gamma \fCenter \Delta$
\DisplayProof
\]
ويعطينا فرض الاستقراء !!a{proof}~$\pi_1''$ لــ$!B, \Gamma \fCenter
\Delta$. تأمل الآن !!{proof} المنتهي بــ\CutCS{}:
\[
\AxiomC{}
\RightLabel{$\pi_2'$}
\Deduce$\Gamma \fCenter \Delta, !B$
\AxiomC{}
\RightLabel{$\pi_1''$}
\Deduce$!B, \Gamma \fCenter \Delta$
\RightLabel{\CutCS}
\BinaryInf$\Gamma \fCenter \Delta$
\DisplayProof
\]
رتبة قطعه أيضًا أصغر من رتبة القطع في~$\pi$، ولذلك يعطينا فرض الاستقراء
!!{proof}~$\pi'$ في \Log{G3c} لــ$\Gamma \fCenter \Delta$.
\item إذا كان $!A \ident \lforall[x][!B(x)]$، فإن $\pi$ ينتهي بــ
\[
\AxiomC{}
\RightLabel{$\pi_1'$}
\Deduce$\Gamma \fCenter \Delta, !B(c)$
\RightLabel{\RightR{\lforall}}
\UnaryInf$\Gamma \fCenter \Delta, \lforall[x][!B(x)]$
\LeftSubproofLabel{$\pi_1$}
\AxiomC{}
\RightLabel{$\pi_2'$}
\Deduce$!B(t), \lforall[x][!B(x)], \Gamma \fCenter \Delta$
\RightLabel{\LeftR{\lforall}}
\UnaryInf$\lforall[x][!B(x)], \Gamma \fCenter \Delta$
\RightSubproofLabel{$\pi_2$}
\RightLabel{\CutCS}
\BinaryInf$\Gamma \fCenter \Delta$
\DisplayProof
\]
وبحسب فرض الاستقراء يمكن حذف \CutCS{} من
\[
\AxiomC{}
\RightLabel{$\pi_1''$}
\Deduce$!B(t), \Gamma \fCenter \Delta, \lforall[x][!B(x)]$
\AxiomC{}
\RightLabel{$\pi_2'$}
\Deduce$!B(t), \lforall[x][!B(x)], \Gamma \fCenter \Delta$
\RightLabel{\CutCS}
\BinaryInf$!B(t), \Gamma \fCenter \Delta$
\DisplayProof
\]
حيث يُستخرج $\pi_1''$ من $\pi_1$ (لا من~$\pi_1'$!) بالإضعاف الحافظ للــ
!!{height}. (لهذا !!{proof} رتبة القطع نفسها، لكن !!{height} قطعه أقل.)
وهذا يعطينا !!a{proof}~$\pi_2''$ لــ$!B(t), \Gamma \fCenter \Delta$.

المتتالية النهائية لــ!!{proof}~$\pi_1'$ هي $\Gamma \Sequent \Delta,
!B(c)$، حيث $c$ المتغير المميز لاستدلال~$\RightR{\lforall}$. ولذلك لا يرد
$c$ في~$!B(x)$ أو $\Gamma$ أو~$\Delta$. نفترض أن !!{proof}~$\pi$ منتظم،
ومن ثم لا يكون $c$ المتغير المميز إلا لاستدلال~$\RightR{\lforall}$ الآتي،
ولا يرد إلا فوقه، أي لا يرد إلا في~$\pi_1'$، وليس متغيرًا مميزًا لاستدلال
في~$\pi_1'$. ولما كان $c$ لا يرد في~$\pi_2$، فلا يمكن أن يرد في~$t$.
وبحسب \olref[seq][qua]{lem:subst-regular} يكون
!!{proof}~$\Subst{\pi_1'}{t}{c}$ هو !!a{proof} لــ$\Gamma \Sequent
\Delta, !B(t)$. ونطبّق الآن فرض الاستقراء على !!{proof}
\[
\AxiomC{}
\RightLabel{$\Subst{\pi_1'}{t}{c}$}
\Deduce$\Gamma \fCenter \Delta, !B(t)$
\AxiomC{}
\RightLabel{$\pi_2''$}
\Deduce$!B(t), \Gamma \fCenter \Delta$
\RightLabel{\CutCS}
\BinaryInf$\Gamma \fCenter \Delta$
\DisplayProof
\]
(ينطبق فرض الاستقراء لأن !!{formula} القطع هي~$!B(t)$، ولذلك تكون رتبة
القطع في !!{proof} أدنى منها في~$\pi$.)
\item نترك الحالة $!A \ident \lexists[x][!B(x)]$ تمرينًا.
\end{enumerate}
\end{proof}

\begin{prob}
تحقق من الحالة الفرعية للحالة~D في برهان \olref{lem:cut-adm-G3c} التي لا
تكون فيها $!A$ رئيسة في الاستدلال الأخير لــ$\pi_1$، المنتهي بــ
$\LeftR{\lif}$. (ملاحظة: يقتضي هذا التمرين جزئيًا أن تصرّح بوضوح بما عليك
إثباته، أي بالصورة التي يكون عليها $\pi$ في هذه الحالة.)
\end{prob}

\begin{prob}
تحقق من الحالة الفرعية للحالة~D في برهان \olref{lem:cut-adm-G3c} التي لا
تكون فيها $!A$ رئيسة في الاستدلال الأخير لــ$\pi_1$، المنتهي بــ
$\LeftR{\lforall}$.
\end{prob}

\begin{prob}
تحقق من الحالة الفرعية للحالة~E في برهان \olref{lem:cut-adm-G3c} التي تكون
فيها $!A$ رئيسة في الاستدلالين الأخيرين لكل من $\pi_1$ و$\pi_2$، ويكون
$!A \ident (!B \land !C)$.
\end{prob}

\begin{prob}
تحقق من الحالة الفرعية للحالة~E في برهان \olref{lem:cut-adm-G3c} التي تكون
فيها $!A$ رئيسة في الاستدلالين الأخيرين لكل من $\pi_1$ و$\pi_2$، ويكون
$!A \ident \lexists[x][!B(x)]$.
\end{prob}

لاحظ أن !!{height} القطع في !!{proof} المنتهي بــ\CutCS{} الذي نطبّق عليه
فرض الاستقراء في الحالة~E قد يكون أكبر من !!{height} القطع في البرهان
الأصلي. فمثلًا، في الحالة التي يكون فيها $!A\ident(!B \lif !C)$ حوّلنا
$\pi$ إلى !!a{proof} يتضمن استدلالي \CutCS{}،
\[
\AxiomC{}
\RightLabel{$\pi_2'$}
\Deduce$\Gamma \fCenter \Delta, !B$
\AxiomC{}
\RightLabel{$\pi_1'$}
\Deduce$!B, \Gamma \fCenter \Delta, !C$
\AxiomC{}
\RightLabel{$\pi_2''$}
\Deduce$!C, \Gamma \fCenter \Delta$
\RightLabel{\CutCS}
\BinaryInf$!B, \Gamma \fCenter \Delta$
\RightSubproofLabel{$\pi_1''$}
\RightLabel{\CutCS}
\BinaryInf$\Gamma \fCenter \Delta$
\DisplayProof
\]
استدلال \CutCS{} العلوي بين $\pi_1'$ و$\pi_2''$ أدنى من~$\pi$ في كل من
رتبة القطع و!!{height} القطع. لكن !!{proof}~$\pi_1''$---وهو نتيجة تطبيق
فرض الاستقراء---لم يعد أدنى من~$\pi$ في !!{height} القطع. غير أن
!!{formula} القطع في استدلال \CutCS{} السفلي هي $!B$، ولذلك تكون
\emph{رتبته} أدنى من رتبة القطع في~$\pi$.

تثبت \olref{lem:cut-adm-G3c} أنه يمكننا حذف استدلال \CutCS{} علوي واحد من
!!a{proof}. ويمكننا إثبات إمكان حذف أي عدد من استدلالات \CutCS{} من
!!a{proof} بتطبيقها تباعًا على !!{proof}s الفرعية العليا المنتهية
بــ\CutCS{}.

\begin{thm}\ollabel{thm:cut-elim-G3c-topmost} يمكن تحويل أي !!{proof}~$\pi$
في $\Log{G3c} + \CutCS$ إلى !!a{proof} في~\Log{G3c}.
\end{thm}

\begin{proof}
  بالاستقراء على عدد استدلالات \CutCS{} في~$\pi$، ولنسمّه~$n$. إذا كان
  $n=0$، فلا شيء نفعله: يكون $\pi$ أصلًا !!a{proof} في~\Log{G3c}. وإلا،
  فتأمل !!{proof} الفرعي~$\pi'$ المنتهي باستدلال \CutCS{} علوي. فهو يحتوي
  على قاعدة \CutCS{} واحدة بوصفها استدلاله الأخير. طبّق
  \olref{lem:cut-adm-G3c} لتحوّله إلى برهان~$\pi''$ خال من \CutCS{}،
  واستبدل !!{proof} الفرعي~$\pi'$ بــ$\pi''$. ينتج !!a{proof} يحتوي
  $n-1$ استدلالًا من نوع \CutCS{}، فينطبق فرض الاستقراء.
\end{proof}

\end{document}
